% Part IX — The Geometry of Ideation  (Slides 201 – 218)
\providecommand{\jugmark}[1]{\textbf{\color{jugeoBlue}#1}}

% ---------------------------------------------------------------------------
% Slide 201: Section divider
% ---------------------------------------------------------------------------
\begin{frame}
\frametitle{}
\vfill
\begin{center}
  {\huge\color{jugeoBlue}\textbf{Part IX}}\\[16pt]
  {\LARGE The Geometry of Ideation}\\[12pt]
  {\large\color{gray}\itshape How JuGeo discovers what to prove next}
\end{center}
\vfill
\end{frame}

% ---------------------------------------------------------------------------
% Slide 202: The Missing Piece — Discovery
% ---------------------------------------------------------------------------
\begin{frame}
\frametitle{The Missing Piece --- Discovery}
\begin{columns}[T]
\column{0.55\textwidth}
So far, JuGeo can:
\begin{itemize}
  \item Model code as a geometric site
  \item Attach evidence with trust levels
  \item Detect and repair failures
  \item Ship proof-carrying Python
\end{itemize}
\vspace{10pt}
But \jugmark{who decides what to verify}?
\begin{itemize}
  \item Which functions need specs?
  \item What invariants might hold?
  \item What ideas are worth pursuing?
\end{itemize}
\column{0.42\textwidth}
\begin{center}
\begin{tikzpicture}[
  box/.style={draw, thick, rounded corners, fill=jugeoLight,
              minimum width=2.8cm, minimum height=0.9cm, font=\small,
              align=center}
]
  \node[box] (v) at (0,0)    {Verification\\Engine};
  \node[box, fill=jugeoGreen!15] (d) at (0,2) {\jugmark{Discovery}\\Engine};
  \node[font=\Large, jugeoBlue] at (0, 1) {?};
  \draw[-{Stealth}, thick, jugeoBlue!50, dashed] (d) -- (v);
\end{tikzpicture}
\end{center}
\end{column}
\end{columns}
\end{frame}

% ---------------------------------------------------------------------------
% Slide 203: Analogy — Running a Research Lab
% ---------------------------------------------------------------------------
\begin{frame}
\frametitle{Analogy: Running a Research Lab}
\begin{center}
\begin{tikzpicture}[
  role/.style={draw, thick, rounded corners, minimum width=2.6cm,
               minimum height=0.85cm, font=\small, align=center, text=white},
  arr/.style={-{Stealth[length=3mm]}, thick, jugeoBlue!60}
]
  \node[role, fill=jugeoOrange]                  (gen) {Generate\\Ideas};
  \node[role, fill=jugeoBlue, right=0.6cm of gen]  (eval) {Evaluate\\Novelty};
  \node[role, fill=jugeoGreen!80, right=0.6cm of eval] (inv) {Invest\\Resources};
  \node[role, fill=jugeoDark, right=0.6cm of inv]  (pub) {Validate \&\\Publish};

  \draw[arr] (gen) -- (eval);
  \draw[arr] (eval) -- (inv);
  \draw[arr] (inv) -- (pub);
\end{tikzpicture}
\end{center}
\vfill
\begin{itemize}
  \item A \textbf{research lab} doesn't just execute experiments blindly
  \item It \textbf{generates} candidate ideas, \textbf{evaluates} their novelty,
        \textbf{invests} limited budget wisely, and \textbf{validates} results
  \item JuGeo's \jugmark{ideation engine} does the same thing for
        code verification and theorem discovery
  \item It treats \textbf{ideas as structured objects} --- not strings, not hunches
\end{itemize}
\end{frame}

% ---------------------------------------------------------------------------
% Slide 204: Ideas as Structured Objects
% ---------------------------------------------------------------------------
\begin{frame}[fragile]
\frametitle{Ideas as Structured Objects}
\begin{center}
\begin{tikzpicture}[
  comp/.style={draw=jugeoBlue, fill=jugeoLight, rounded corners=3pt,
               minimum width=26mm, minimum height=8mm, font=\small\bfseries,
               line width=0.7pt, align=center},
  arr/.style={-{Stealth[length=4pt]}, jugeoBlue!50, line width=0.5pt}
]
  \node[comp] (prop) at (0, 1.5)   {Proposition};
  \node[comp] (nov)  at (3.5, 1.5) {Novelty Score};
  \node[comp] (feas) at (7, 1.5)   {Feasibility};
  \node[comp] (dep)  at (0, 0)     {Dependencies};
  \node[comp] (cost) at (3.5, 0)   {Cost Estimate};
  \node[comp] (trust) at (7, 0)    {Trust Status};

  \draw[arr] (prop) -- (nov);
  \draw[arr] (nov) -- (feas);
  \draw[arr] (prop) -- (dep);
  \draw[arr] (dep) -- (cost);
  \draw[arr] (feas) -- (trust);
  \draw[arr] (cost) -- (trust);
\end{tikzpicture}
\end{center}
\vfill
An \jugmark{Idea} in JuGeo is not just a string --- it's a typed object:
\begin{itemize}
  \item \textbf{Proposition}: What is being claimed?
  \item \textbf{Novelty}: How different is this from what we already know?
  \item \textbf{Feasibility}: Can we actually prove or validate it?
  \item \textbf{Dependencies}: What must be true first?
  \item \textbf{Cost}: How much effort would it take?
  \item \textbf{Trust status}: Proposed? Partially validated? Fully proven?
\end{itemize}
\end{frame}

% ---------------------------------------------------------------------------
% Slide 205: Novelty — Measuring Distance from the Known
% ---------------------------------------------------------------------------
\begin{frame}
\frametitle{Novelty --- Distance from the Known}
\begin{columns}[T]
\column{0.52\textwidth}
\begin{block}{The Question}
Is this idea \textbf{actually new}, or just a restatement of something
we already know?
\end{block}
\vspace{6pt}
JuGeo measures novelty as \textbf{semantic distance}:
\begin{itemize}
  \item Maintain a \jugmark{theorem portfolio} of known results
  \item Score each candidate against the portfolio
  \item High novelty = far from existing knowledge
  \item But novelty alone isn't enough ---\\
        the idea must also be \textbf{useful}
\end{itemize}
\column{0.45\textwidth}
\begin{center}
\begin{tikzpicture}[scale=0.85]
  % Known territory
  \fill[jugeoLight] (0,0) circle (1.8);
  \node[font=\scriptsize\bfseries, jugeoBlue] at (0,0) {Known};
  % Ideas
  \fill[jugeoGreen!80] (0.8, 0.3) circle (3pt);
  \node[font=\tiny, right] at (0.95, 0.3) {low novelty};
  \fill[jugeoOrange] (2.4, 1.2) circle (3pt);
  \node[font=\tiny, right] at (2.55, 1.2) {high novelty};
  \fill[jugeoRed] (3.5, 2.5) circle (3pt);
  \node[font=\tiny, right] at (3.65, 2.5) {too far (infeasible)};
  % Sweet spot
  \draw[jugeoGreen, thick, dashed] (0,0) circle (2.8);
  \node[font=\tiny\itshape, jugeoGreen!80!black, rotate=30] at (1.5, 2.3)
    {sweet spot};
\end{tikzpicture}
\end{center}
\end{column}
\end{columns}
\end{frame}

% ---------------------------------------------------------------------------
% Slide 206: Regimes — Different Modes of Discovery
% ---------------------------------------------------------------------------
\begin{frame}
\frametitle{Regimes --- Different Modes of Discovery}
\begin{center}
\begin{tikzpicture}[
  reg/.style={draw, thick, rounded corners, minimum width=2.8cm,
              minimum height=1.3cm, font=\small, align=center}
]
  \node[reg, fill=jugeoBlue!15]
    (exp) at (0, 0) {\textbf{Exploration}\\Find new territory};
  \node[reg, fill=jugeoGreen!15]
    (ref) at (4.5, 0) {\textbf{Refinement}\\Deepen existing results};
  \node[reg, fill=jugeoOrange!15]
    (bri) at (9, 0) {\textbf{Bridging}\\Connect two domains};

  \draw[-{Stealth}, thick, jugeoBlue!40, bend left=20]
    (exp) to node[above, font=\tiny]{saturated} (ref);
  \draw[-{Stealth}, thick, jugeoGreen!40, bend left=20]
    (ref) to node[above, font=\tiny]{stuck} (bri);
  \draw[-{Stealth}, thick, jugeoOrange!40, bend left=20]
    (bri) to node[below, font=\tiny]{new territory} (exp);
\end{tikzpicture}
\end{center}
\vfill
\begin{itemize}
  \item Like a GPS with multiple route modes, JuGeo picks the right
        \jugmark{discovery regime} for the current situation
  \item \textbf{Exploration}: cast a wide net when the search space is open
  \item \textbf{Refinement}: go deeper when existing results can be strengthened
  \item \textbf{Bridging}: connect two domains when each is stuck alone
  \item Regimes transition automatically based on progress and diminishing returns
\end{itemize}
\end{frame}

% ---------------------------------------------------------------------------
% Slide 207: Analogy Transport — Moving Ideas Between Domains
% ---------------------------------------------------------------------------
\begin{frame}
\frametitle{Analogy Transport}
\begin{center}
\begin{tikzpicture}[
  dom/.style={draw, thick, rounded corners, minimum width=2.4cm,
              minimum height=1.5cm, font=\small, align=center},
  arr/.style={-{Stealth[length=4mm]}, thick, jugeoGreen}
]
  \node[dom, fill=jugeoBlue!10]
    (a) at (0, 0) {\textbf{Sorting}\\algorithms\\[2pt]\scriptsize ``merge preserves\\order''};
  \node[dom, fill=jugeoOrange!10]
    (b) at (6, 0) {\textbf{Database}\\merge joins\\[2pt]\scriptsize ``merge preserves\\keys''};

  \draw[arr] (a) -- node[above, font=\small\bfseries]{Analogy Transport}
                    node[below, font=\scriptsize\itshape]{structure-preserving map} (b);
\end{tikzpicture}
\end{center}
\vfill
\begin{itemize}
  \item A theorem about sorting can sometimes be \jugmark{transported}
        to a structurally similar domain
  \item ``Merge preserves order'' $\longrightarrow$ ``Merge join preserves key uniqueness''
  \item JuGeo finds these structural similarities and transports ideas
        across domains --- \textbf{with trust attenuation}
  \item The transported idea starts at \textsc{proposed} trust\\
        (never silently promoted to \textsc{verified})
  \item This is how discoveries in one codebase can accelerate
        verification of another
\end{itemize}
\end{frame}

% ---------------------------------------------------------------------------
% Slide 208: The Discovery Pipeline
% ---------------------------------------------------------------------------
\begin{frame}
\frametitle{The Discovery Pipeline}
\begin{center}
\begin{tikzpicture}[
  stage/.style={draw, thick, rounded corners, minimum width=2.2cm,
                minimum height=0.85cm, font=\small\bfseries, text=white, align=center},
  arr/.style={-{Stealth[length=3mm]}, thick, jugeoBlue!60}
]
  \node[stage, fill=jugeoOrange]                         (gen) {Generate};
  \node[stage, fill=jugeoBlue,   right=0.4cm of gen]     (nov) {Score\\Novelty};
  \node[stage, fill=jugeoBlue!80, right=0.4cm of nov]    (cls) {Classify\\Kind};
  \node[stage, fill=jugeoGreen!80, right=0.4cm of cls]   (val) {Validate};
  \node[stage, fill=jugeoDark,    right=0.4cm of val]    (pro) {Promote};

  \draw[arr] (gen) -- (nov);
  \draw[arr] (nov) -- (cls);
  \draw[arr] (cls) -- (val);
  \draw[arr] (val) -- (pro);
\end{tikzpicture}
\end{center}
\vspace{8pt}
\begin{description}[\textbf{Promote}~~]
  \item[\textbf{Generate}] Produce candidate ideas from code analysis,
        analogy, or AI suggestion
  \item[\textbf{Score}] Measure novelty against the known portfolio
  \item[\textbf{Classify}] What kind of idea is it? Invariant? Equivalence?
        Bug pattern? New concept?
  \item[\textbf{Validate}] Test the idea --- run experiments, check with the solver,
        attempt a proof
  \item[\textbf{Promote}] If validated, inject into JuGeo's judgment sheaf
        as a coordinate with trust and evidence
\end{description}
\end{frame}

% ---------------------------------------------------------------------------
% Slide 209: Theorem Economics — Investing Wisely
% ---------------------------------------------------------------------------
\begin{frame}
\frametitle{Theorem Economics --- Investing Wisely}
\begin{columns}[T]
\column{0.48\textwidth}
\begin{block}{The Problem}
We have \textbf{limited budget} (time, compute, human attention)
and \textbf{many candidate ideas}. Which should we pursue?
\end{block}
\vspace{6pt}
JuGeo treats this as an \jugmark{investment problem}:
\begin{itemize}
  \item Each idea has an expected \textbf{yield}\\
        (how much verified coverage it adds)
  \item Each idea has a \textbf{cost}\\
        (solver time, human effort, risk of failure)
  \item \textbf{Diminishing returns} mean switching topics is sometimes
        better than going deeper
\end{itemize}
\column{0.48\textwidth}
\begin{center}
\begin{tikzpicture}[scale=0.85]
  \draw[-{Stealth}, thick] (0,0) -- (4.5,0) node[right, font=\scriptsize]{effort};
  \draw[-{Stealth}, thick] (0,0) -- (0,3.5) node[above, font=\scriptsize]{value};
  % Idea A — high initial return, saturates
  \draw[jugeoGreen, thick] plot[smooth]
    coordinates {(0,0) (0.5,1.5) (1,2.2) (1.5,2.6) (2,2.8) (3,2.9) (4,2.95)};
  \node[font=\scriptsize, jugeoGreen] at (4.2, 3.1) {Idea A};
  % Idea B — slow start, keeps growing
  \draw[jugeoBlue, thick] plot[smooth]
    coordinates {(0,0) (0.5,0.3) (1,0.8) (1.5,1.5) (2,2.3) (3,3.1) (4,3.4)};
  \node[font=\scriptsize, jugeoBlue] at (4.2, 3.5) {Idea B};
  % Switch point
  \draw[jugeoRed, dashed, thick] (1.7, 0) -- (1.7, 3.2);
  \node[font=\tiny, jugeoRed, align=center] at (1.7, -0.4) {switch\\point};
\end{tikzpicture}
\end{center}
\end{column}
\end{columns}
\end{frame}

% ---------------------------------------------------------------------------
% Slide 210: Theorem Ecologies — Ideas That Help Each Other
% ---------------------------------------------------------------------------
\begin{frame}
\frametitle{Theorem Ecologies --- Ideas That Help Each Other}
\begin{center}
\begin{tikzpicture}[
  thm/.style={draw, thick, circle, minimum size=10mm, font=\scriptsize\bfseries,
              fill=jugeoLight, inner sep=1pt},
  dep/.style={-{Stealth[length=3pt]}, thick, jugeoBlue!50}
]
  \node[thm, fill=jugeoGreen!20] (a) at (0, 0)    {A};
  \node[thm] (b) at (2, 1.2)  {B};
  \node[thm] (c) at (4, 0)    {C};
  \node[thm, fill=jugeoGreen!20] (d) at (2, -1.2) {D};
  \node[thm] (e) at (5.5, 1)  {E};
  \node[thm, fill=jugeoOrange!20] (f) at (5.5, -1) {F};

  \draw[dep] (a) -- (b);
  \draw[dep] (a) -- (d);
  \draw[dep] (b) -- (c);
  \draw[dep] (d) -- (c);
  \draw[dep] (c) -- (e);
  \draw[dep] (c) -- (f);

  \node[font=\tiny\itshape, gray, below=2pt of a] {proven};
  \node[font=\tiny\itshape, gray, below=2pt of d] {proven};
  \node[font=\tiny\itshape, jugeoOrange, below=2pt of f] {high yield};
\end{tikzpicture}
\end{center}
\vfill
\begin{itemize}
  \item Theorems form an \jugmark{ecology}: proving A makes B easier,
        which unlocks C, which enables F
  \item JuGeo tracks these \textbf{compounding effects}
  \item A cheap theorem that unlocks expensive ones is worth more than
        its individual yield
  \item The economics layer optimizes for \textbf{portfolio value},
        not just individual theorems
\end{itemize}
\end{frame}

% ---------------------------------------------------------------------------
% Slide 211: Semantic Futures — Where Could We Go?
% ---------------------------------------------------------------------------
\begin{frame}
\frametitle{Semantic Futures --- Where Could We Go?}
\begin{center}
\begin{tikzpicture}[
  state/.style={draw, thick, rounded corners, minimum width=2cm,
                minimum height=0.7cm, font=\scriptsize, align=center},
  now/.style={state, fill=jugeoGreen!20, line width=1.5pt, draw=jugeoGreen},
  fut/.style={state, fill=jugeoLight},
  arr/.style={-{Stealth[length=3pt]}, thick, jugeoBlue!40}
]
  \node[now] (now) at (0, 0) {\textbf{Current State}\\5 theorems proven};
  \node[fut] (f1) at (-3.5, 1.8) {State A\\7 theorems};
  \node[fut] (f2) at (0, 2.2) {State B\\8 theorems};
  \node[fut] (f3) at (3.5, 1.8) {State C\\6 theorems};
  \node[fut, fill=jugeoGreen!10, draw=jugeoGreen] (f4) at (0, 3.8)
    {State D\\12 theorems};

  \draw[arr] (now) -- (f1);
  \draw[arr] (now) -- (f2);
  \draw[arr] (now) -- (f3);
  \draw[arr] (f1) -- (f4);
  \draw[arr] (f2) -- (f4);
  \draw[arr] (f3) -- (f4);
\end{tikzpicture}
\end{center}
\vfill
\begin{itemize}
  \item JuGeo models the space of \jugmark{reachable future states}
  \item Each future state represents a set of proven theorems
  \item The ideation engine asks: ``which futures are most valuable
        and most reachable?''
  \item Uses beam search, diversity constraints, and budget allocation
        to navigate toward high-value futures
\end{itemize}
\end{frame}

% ---------------------------------------------------------------------------
% Slide 212: The Synthesis Frontier — Cross-Domain Invention
% ---------------------------------------------------------------------------
\begin{frame}
\frametitle{The Synthesis Frontier}
\begin{center}
\begin{tikzpicture}[
  field/.style={draw, thick, rounded corners, minimum width=2cm,
                minimum height=0.8cm, font=\small, align=center, fill=jugeoLight},
  arr/.style={-{Stealth[length=3pt]}, thick, jugeoOrange}
]
  \node[field] (a) at (-2.5, 1)  {Algebra};
  \node[field] (b) at (2.5, 1)   {Networking};
  \node[field, fill=jugeoGreen!15, draw=jugeoGreen, line width=1.2pt]
    (synth) at (0, -0.8)
    {\textbf{New Proposition}};

  \draw[arr] (a) -- (synth);
  \draw[arr] (b) -- (synth);

  \node[font=\scriptsize\itshape, gray] at (0, 0.5) {structural overlap};
\end{tikzpicture}
\end{center}
\vfill
The \jugmark{synthesis frontier} combines ideas from different fields:
\begin{itemize}
  \item JuGeo maintains a catalog of \textbf{128 knowledge fields}
  \item It finds \textbf{structural overlaps} between field pairs
  \item Generates cross-domain propositions via \textbf{metaphor and pattern matching}
  \item Runs a \textbf{tournament} to select the most promising candidates
  \item Winners enter the judgment sheaf at \textsc{proposal} trust
\end{itemize}
\end{frame}

% ---------------------------------------------------------------------------
% Slide 213: From Idea to Judgment — The Bridge
% ---------------------------------------------------------------------------
\begin{frame}
\frametitle{From Idea to Judgment}
\begin{center}
\begin{tikzpicture}[
  box/.style={draw, thick, rounded corners, minimum width=2.8cm,
              minimum height=1cm, font=\small, align=center},
  arr/.style={-{Stealth[length=4mm]}, thick}
]
  \node[box, fill=jugeoOrange!15]              (idea) {Idea\\(unvalidated)};
  \node[box, fill=jugeoBlue!15, right=1.5cm of idea] (val)  {Validated\\Candidate};
  \node[box, fill=jugeoGreen!15, right=1.5cm of val]  (judg) {Judgment\\$(c, \varphi, A, E, O, B, T, \Pi)$};

  \draw[arr, jugeoOrange] (idea) -- node[above, font=\scriptsize]{experiment} (val);
  \draw[arr, jugeoGreen]  (val)  -- node[above, font=\scriptsize]{inject} (judg);
\end{tikzpicture}
\end{center}
\vfill
\begin{itemize}
  \item Every validated idea becomes a \jugmark{judgment 8-tuple}
  \item Gets a coordinate in the site, a proposition, initial evidence,
        obligations, and \textsc{proposal}-level trust
  \item From there, the normal verification pipeline takes over:
        descent, evidence routing, trust promotion
  \item \jugmark{Trust discipline is preserved}: an idea from AI starts low
        and must earn its way up
\end{itemize}
\end{frame}

% ---------------------------------------------------------------------------
% Slide 214: Theory Navigation — Exploring the Map
% ---------------------------------------------------------------------------
\begin{frame}
\frametitle{Theory Navigation --- Exploring the Map}
\begin{center}
\begin{tikzpicture}[
  node/.style={draw, thick, circle, minimum size=6mm, inner sep=0pt,
               font=\tiny\bfseries},
  visited/.style={node, fill=jugeoGreen!30},
  current/.style={node, fill=jugeoOrange!40, line width=1.5pt},
  unknown/.style={node, fill=gray!15},
  edge/.style={thick, gray!40},
  path/.style={thick, jugeoGreen!70}
]
  \node[visited] (a) at (0, 0)    {A};
  \node[visited] (b) at (1.5, 1)  {B};
  \node[current] (c) at (3, 0.5)  {C};
  \node[unknown] (d) at (4.5, 1.5){D};
  \node[unknown] (e) at (4.5, -0.5){E};
  \node[visited] (f) at (1, -1.2) {F};
  \node[unknown] (g) at (6, 0.5)  {G};

  \draw[edge] (a)--(b);
  \draw[path] (b)--(c);
  \draw[edge] (c)--(d);
  \draw[edge] (c)--(e);
  \draw[edge] (a)--(f);
  \draw[edge] (d)--(g);
  \draw[edge] (e)--(g);
\end{tikzpicture}
\end{center}
\vfill
\begin{itemize}
  \item The entire space of ideas forms a \jugmark{navigable theory graph}
  \item \textcolor{jugeoGreen!70}{Green} = explored,
        \textcolor{jugeoOrange}{orange} = current position,
        \textcolor{gray}{gray} = unexplored
  \item JuGeo tracks maturity, coverage, and frontier novelty
  \item Navigation strategies: depth-first (specialize), breadth-first
        (survey), purpose-directed (find specific things)
  \item This is how JuGeo plans an entire \textbf{research programme},
        not just individual proofs
\end{itemize}
\end{frame}

% ---------------------------------------------------------------------------
% Slide 215: Kind Discovery — Finding New Types of Things
% ---------------------------------------------------------------------------
\begin{frame}
\frametitle{Kind Discovery --- Finding New Types of Things}
\begin{columns}[T]
\column{0.55\textwidth}
Sometimes the \textbf{existing categories} aren't enough.
\vspace{8pt}
\begin{exampleblock}{Example}
You've been verifying sorting algorithms. Then you encounter a
\textbf{streaming sort} that doesn't fit any existing pattern.
\end{exampleblock}
\vspace{8pt}
JuGeo can detect when it encounters something that doesn't fit
existing kinds and \jugmark{bootstrap a new regime}:
\begin{itemize}
  \item Detect anomalous obstruction patterns
  \item Identify new structural regularities
  \item Create a new ``chart'' in the theory atlas
  \item Populate it with initial theorems
\end{itemize}
\column{0.42\textwidth}
\begin{center}
\begin{tikzpicture}[
  chart/.style={draw, thick, rounded corners, minimum width=20mm,
                minimum height=14mm, font=\scriptsize, align=center}
]
  \node[chart, fill=jugeoBlue!10] (c1) at (0, 1.2) {Sorting\\algorithms};
  \node[chart, fill=jugeoGreen!10] (c2) at (0, -0.8) {Search\\algorithms};
  \node[chart, fill=jugeoOrange!10, dashed, draw=jugeoOrange, line width=1.2pt]
    (c3) at (3, 0.2) {\jugmark{NEW:}\\Streaming\\algorithms};

  \draw[-{Stealth}, jugeoOrange, thick, dashed] (1.4, 0.6) -- (c3.west);
  \node[font=\tiny\itshape, jugeoOrange, rotate=15] at (1.7, 1) {bootstrap};
\end{tikzpicture}
\end{center}
\end{column}
\end{columns}
\end{frame}

% ---------------------------------------------------------------------------
% Slide 216: Federated Discovery — Multiple Discovery Sources
% ---------------------------------------------------------------------------
\begin{frame}
\frametitle{Federated Discovery}
\begin{center}
\begin{tikzpicture}[
  src/.style={draw, thick, rounded corners, minimum width=2.2cm,
              minimum height=0.8cm, font=\small, align=center, fill=jugeoLight},
  hub/.style={draw, thick, rounded corners=8pt, minimum width=3cm,
              minimum height=1cm, font=\small\bfseries, align=center,
              fill=jugeoGreen!15, draw=jugeoGreen, line width=1.2pt},
  arr/.style={-{Stealth[length=3mm]}, thick, jugeoBlue!50}
]
  \node[src] (s1) at (-3.5, 1.3) {Code Analysis};
  \node[src] (s2) at (-3.5, 0)   {AI Suggestions};
  \node[src] (s3) at (-3.5,-1.3) {Human Input};
  \node[src] (s4) at (3.5, 1.3)  {Analogy Engine};
  \node[src] (s5) at (3.5, 0)    {Other Projects};
  \node[src] (s6) at (3.5, -1.3) {Literature};

  \node[hub] (fed) at (0, 0) {Federation\\Hub};

  \draw[arr] (s1) -- (fed);
  \draw[arr] (s2) -- (fed);
  \draw[arr] (s3) -- (fed);
  \draw[arr] (s4) -- (fed);
  \draw[arr] (s5) -- (fed);
  \draw[arr] (s6) -- (fed);
\end{tikzpicture}
\end{center}
\vfill
\begin{itemize}
  \item Ideas come from \jugmark{multiple sources} with different trust levels
  \item The \textbf{federation hub} reconciles, deduplicates, and ranks them
  \item Trust is preserved: AI suggestions $\neq$ formal analysis
  \item Conflicts between sources are resolved by consensus or escalation
\end{itemize}
\end{frame}

% ---------------------------------------------------------------------------
% Slide 217: Why Ideation Matters
% ---------------------------------------------------------------------------
\begin{frame}
\frametitle{Why Ideation Matters}
\begin{columns}[T]
\column{0.48\textwidth}
\begin{alertblock}{Without Ideation}
  \begin{itemize}
    \item Human must decide what to verify
    \item No systematic exploration
    \item Missed opportunities
    \item Verification is \textbf{reactive}
  \end{itemize}
\end{alertblock}
\column{0.48\textwidth}
\begin{exampleblock}{With Ideation}
  \begin{itemize}
    \item System discovers what to verify
    \item Novelty-guided exploration
    \item Cross-domain transport of insights
    \item Verification is \textbf{proactive}
  \end{itemize}
\end{exampleblock}
\end{columns}
\vfill
\begin{center}
\begin{tikzpicture}
  \node[fill=jugeoLight, rounded corners, inner sep=10pt, font=\small,
        text width=10cm, align=center]
    {\itshape JuGeo doesn't just check your code ---
     it \jugmark{discovers things about your code} you didn't know to ask.};
\end{tikzpicture}
\end{center}
\end{frame}

% ---------------------------------------------------------------------------
% Slide 218: Ideation Summary
% ---------------------------------------------------------------------------
\begin{frame}
\frametitle{Ideation Summary}
\begin{center}
\small
\begin{tabular}{@{} r l @{}}
  \toprule
  \textbf{Component} & \textbf{What It Does} \\
  \midrule
  Ideas              & Structured proposals with novelty, cost, and trust \\
  Novelty Search     & Purpose-conditioned distance from known theorems \\
  Regimes            & Exploration / refinement / bridging modes \\
  Analogy Transport  & Move ideas across domains, preserving structure \\
  Discovery Pipeline & Generate $\to$ score $\to$ classify $\to$ validate $\to$ promote \\
  Theorem Economics  & Budget-optimal investment across candidate ideas \\
  Theorem Ecologies  & Track compounding effects between theorems \\
  Semantic Futures   & Model and navigate reachable future states \\
  Synthesis Frontier & Cross-domain invention via 128-field catalog \\
  Kind Discovery     & Bootstrap new conceptual regimes when needed \\
  Federation         & Reconcile ideas from multiple sources with trust \\
  \bottomrule
\end{tabular}
\end{center}
\vfill
\begin{center}
\large Next: \jugmark{The Geometry of Orchestration} --- coordinating the whole machine
\end{center}
\end{frame}
